\section{物理一致的生成式逆向设计}

\subsection{仿真保持物理权威的来源}

起点是与角度分辨透射响应配对的自由形态二值结构的仿真基础数据集。这一选择直接源于问题的物理。对于傅里叶算子超表面，最终重要的量是响应场的形状，而该场必须通过麦克斯韦一致的电磁评估来判断。在本项目中，RCWA用于在固定物理堆栈和材料系统上对波长--角度网格上的每个结构进行标注。计算负担是实质性的，但将目标简化为较弱的标量代理会抹除区分真正算子与局部充分器件的结构本身 \cite{moharam1981,moharam1995}。

这种监督的计算负担正是使自由形态逆向设计困难的原因。每个候选原则上必须在有限角度和光谱孔径上进行评估，且结果响应必须用算子感知标准而非单一透射数字来判断。因此，这里采用的设计循环不试图用习得模型替代物理。仿真锚定程序的两端：它在训练期间提供监督信号，并在候选生成后提供最终决策规则。

\subsection{前向替代模型解决筛选负担，而非物理闭合}

对每个生成候选进行密集RCWA评估足够昂贵，以至于某些习得近似在实际中变得有用。前向替代模型起这一作用。它被训练来近似从结构到响应场的映射，用于通过将注意力集中在更可能满足目标的候选上来减少筛选负担。因此，替代模型最好被理解为求解器分配工具：它告诉我们全波努力值得花费在哪里，但它不决定什么算作有效器件。

这种区别很重要，因为替代模型保真度和物理保真度不可互换。习得的预测器可能恢复响应景观中的广泛趋势，同时错过决定目标曲率、边缘滚降或偏振泄漏是否可接受的局部结构。对于算子超表面，这种差异往往是真实算子行为与名义条件附近视觉合理近似之间的差异。当前代码库直接反映这一原则：前向模型用物理空间重建损失训练，用于指导候选集中，而最终接受推迟到RCWA评估。替代模型减少电磁成本，但只是部分地和有意地。

\subsection{为什么逆向阶段必须是生成式的}

从响应场到几何的逆图本质上是多模态的。即使对于高度结构化的目标，也很少有一个二值布局 deserves 被称为唯一正确解。随着几何类变得更自由和目标描述变得更丰富，这种模糊性变得更加明显。在这种情况下，确定性回归器倾向于返回几个潜在设计族的平均，通常在物理验证后产生较弱的最佳情况行为。生成模型在这里的吸引力在于更窄的原因：它们保留几种不同结构可能满足相同光学约束的可能性。

扩散模型非常适合这一角色，因为它们支持对复杂高维分布的条件采样，并已在光子逆向设计中显示出前景 \cite{liu2018,so2019,ma2021,tanriover2023,zhang2023,zhang2024,hen2026}。这里条件输入是完整的目标响应场而非简化描述符。因此，模型被要求从定义在角度、波长和偏振上的连续光学规范推断几何。那是响应空间表述和生成视角最自然相遇的地方：目标越丰富，单一确定性逆像应足够的说法就越不可信。

\subsection{物理一致性通过闭环验证强制执行}

最后一步是生成后物理验证和重排序。生成候选由电磁求解器重新评估，任务特定分数从验证响应而非仅从替代预测计算。这不是在机器学习模型已经"解决"设计问题后添加的装饰性修正。它是将候选生成转变为物理可信逆向设计程序的步骤。

重排序的需求在项目输出本身中很明显。候选集通常包含具有相似名义分数但在边缘区域、角度肩部、光谱漂移或无源偏振泄漏上行为显著不同的结构。因此，排序脚本分离形状一致性、中心抑制、边缘行为、带宽保持和通道泄漏，而非将所有内容压缩成单一标量。物理一致的设计循环必须足够长时间保持候选多样性以使这种验证比较重要。这就是为什么当前方法最好被理解为\emph{物理一致生成式逆向设计}：候选生成探索可行集，替代模型减少筛选负担，电磁求解器闭合循环。
